% =============================================================================
% Section: Intellectual Context
% =============================================================================

\section{Intellectual Context}
\label{sec:related}

We situate the \Opoch{} framework within five intellectual lineages,
identifying precise points of contact and divergence.

% -------------------------------------------------------------------------
\subsection{Constructive Foundations}
\label{sec:related:constructive}

Martin-L\"{o}f's intuitionistic type theory \citep{martinlof1984} requires
explicit constructions for existence claims; Bishop's constructive
analysis \citep{bishop1967} demonstrated that large portions of classical
mathematics survive under this constraint; Kleene's
realizability \citep{kleene1952} formalized truth as recursive function
existence.  The Curry--Howard correspondence identifies proofs with
programs, connecting logic and computation.  A0$^*$ distills the constructive
programme further: not just ``existence requires a witness'' but
``\emph{distinction} requires a witness.''  Where Martin-L\"{o}f begins
with a rich type structure (universes, dependent types), we begin with
$\noth$ and derive structure.  The self-hosting property (Step~19) is a
Curry--Howard instance: the derivation's proof corresponds to a program
that verifies itself.

% -------------------------------------------------------------------------
\subsection{Computation and Information}
\label{sec:related:computation}

Solomonoff's universal induction \citep{solomonoff1964} and Kolmogorov
complexity \citep{kolmogorov1963,livitanyi2008} define a universal prior
over hypotheses weighted by description length.  Wheeler's ``It from
Bit'' \citep{wheeler1990} proposed that physical quantities derive meaning
from information-theoretic questions.  Zuse \citep{zuse1969} and
Fredkin \citep{fredkin2003} modeled the universe as discrete computation.
Deutsch and Marletto's constructor theory \citep{deutsch2015} reformulates
physics in terms of which transformations can and cannot be performed.
Our carrier $\Carrier$ is precisely the set on which Kolmogorov complexity
is defined, and the prefix-free coding at Step~2 satisfies the same Kraft
inequality \citep{kraft1949}.  The key divergence from Solomonoff is
philosophical: we use $\UTM$ as a \emph{verifier}, not a prior for
probabilistic prediction; the trit $\{\UNIQUE, \OMEGA, \UNSAT\}$ replaces
probability.  Unlike digital physics, our derivation does not assume
discrete computation---Step~3 derives the need for a computational
substrate from A0$^*$; the Church--Turing thesis corroborates the forced class.
Both constructor theory and \Opoch{} ask ``what tasks are possible?'' rather
than ``what trajectories occur?''---but our derivation forces the test
algebra $\Del^*$ from A0$^*$ rather than taking possible transformations as
primitive.

% -------------------------------------------------------------------------
\subsection{Verification and Certification}
\label{sec:related:verification}

Necula's proof-carrying code \citep{necula1997} attaches machine-checkable
proofs to executables; Goldwasser, Micali, and
Rackoff \citep{goldwasser1989} established interactive proof systems;
Clarke et~al.\ \citep{clarke2018} provide the reference for model checking.
McConnell et~al.\ \citep{mcconnell2011} formalized certifying algorithms
that produce efficiently checkable certificates alongside results.
Belnap's four-valued logic \citep{belnap1977} extends classical logic with
``unknown'' and ``contradictory'' values.  The \Opoch{} receipt system is
proof-carrying code applied to reasoning traces: every execution produces a
SHA-256 hash chain that is replay-verifiable.  Unlike interactive proofs,
verification is non-interactive; unlike model checking, the Bellman search
prunes the state space dynamically.  The exactness gate (Step~14) implements
McConnell's certificate-first pattern.  Our trit differs from Belnap's four
values: we exclude ``contradictory'' (a deterministic verifier cannot produce
genuine contradictions), and our $\OMEGA$ is constructive---it returns the
minimal separator test, not a static ``unknown'' label.

% -------------------------------------------------------------------------
\subsection{Geometry, Topology, and Automata}
\label{sec:related:geometry}

Scott's domain theory \citep{scott1970} models computation via continuous
lattices; Smyth \citep{smyth1983} and Abramsky \citep{abramsky1994} connect
topological structure to logical systems.  The Myhill--Nerode
theorem \citep{myhill1957,nerode1958} identifies strings with identical
futures, yielding minimal automata \citep{hopcroft1979}; Milner's
bisimulation \citep{milner1989} extends this to concurrent systems.
Amari's information geometry \citep{amari2016} treats statistical models
as Riemannian manifolds with the Fisher metric measuring
distinguishability \citep{rao1945}.  Our observable opens $\ObsOpen$
(Step~7) are a finite restriction of the Scott topology to testable
predicates.  The Myhill--Nerode congruence $\MNcong$ (Step~15) applies
futures equivalence to the system's state space rather than string classes,
with the test algebra $\Del^*$ as alphabet.  The separator geometry
(Step~16) constructs a metric $\SepDist$ analogous to Fisher-Rao distance
but defined over deterministic truth classes and measured in test cost
rather than information divergence.

% -------------------------------------------------------------------------
\subsection{Decision Theory and Applications}
\label{sec:related:applications}

Bellman's principle of optimality \citep{bellman1957} and Wald's minimax
decision theory \citep{wald1950} provide the mathematical foundation for
the framework's test-selection rule: at each step, select the test minimizing
worst-case cost-to-go.  Chain-of-thought \citep{wei2022} and
tree-of-thoughts \citep{yao2023} improve LLM reasoning by eliciting
intermediate steps, but every step remains a sample from a learned
distribution.  The \Opoch{} test selection (Step~18) is a direct application
of the Bellman equation in test space rather than action space, with
Wald's minimax framing justifying the worst-case optimization.

% -------------------------------------------------------------------------
\subsection{Self-Reference}
\label{sec:related:selfref}

The Kleene recursion theorem guarantees that every total recursive function
can compute its own index; metacircular interpreters and reflective towers
formalize systems reasoning about their own structure.  Our self-hosting
property $\SelfHost = \mathcal{S}$ (Step~19) is the verification analog:
the system verifies its own derivation, producing $\UNIQUE$ with receipts
for each step.  This is stronger than metacircular interpretation (which
merely executes) and connects to proof-carrying code (each step's receipt
is a proof that the step is forced).  Mechanizing the derivation in a proof
assistant (Lean, Coq) would yield a certified executable system via the
Curry--Howard correspondence.
